
\newtoggle{withbonus} \toggletrue{withbonus} % set true to show bonus points in grade table

% setting underlying course name (Grundlagenfach Mathematik, §-Unterricht, \ldots)
\newcommand{\mycourse}{Ergänzungsfach Informatik}
% name of the class
\newcommand{\myclass}{4. Klassen}
% set date
\newcommand{\mydate}{15. Januar 2025}

\title{\textbf{Komplexitätstheorie}}
\author{Thomas Graf\\
	\mycourse, \myclass}
\date{\mydate}
\maketitle
\thispagestyle{headandfoot}

\begin{center}
	\fbox{\parbox{140mm}{\centering
			\begin{itemize}
				\item Dauer der Prüfung: $60$ Minuten
				      %\item \textbf{Zeige in jeder Aufgabe unbedingt Deine Schritte!}
				\item Erlaubte Hilfsmittel:
				      \begin{itemize}
					      \item Schreibzeug
					            % \item Taschenrechner
				      \end{itemize}
			\end{itemize}}}
\end{center}
\vspace{10mm}

\begin{center}
	Vorname: \rule{45mm}{0.8pt}\qquad Nachname: \rule{45mm}{0.8pt}
\end{center}
\vspace{20mm}

\begin{center}
	\iftoggle{withbonus} {
		\combinedgradetable[h][questions]
	} {
		\gradetable[h][questions]
	}
\end{center}
\vspace{15mm}
\begin{center}
	Note:
	\vspace{10mm}

	\rule{8mm}{1.2pt}
\end{center}
% \thispagestyle{empty}
\clearpage


\begin{questions}

	\question[2]

	Welche der folgenden Aussagen sind korrekt?
	\begin{checkboxes}
		\correctchoice $n - 4n^2 \in O(n^3)$ für $n\to\infty$
		\begin{solution}
			richtig, $(n - 4n^2)/(n^3) = 1/(n^2) - 4/n$ ist eine beschränkte Folge
		\end{solution}
		\choice $n - 4n^2 \in O(n)$ für $n\to\infty$
		\begin{solution}
			falsch, $(n - 4n^2)/n = 1-4n$ ist eine unbeschränkte Folge
		\end{solution}
		\correctchoice $n - 4n^2 \in O(n^2)$ für $n\to\infty$
		\begin{solution}
			richtig, $(n - 4n^2)/(n^2) = 1/n - 4$ ist eine beschränkte Folge
		\end{solution}
	\end{checkboxes}


	\question

	Es seien $(a_n)_n$ und $(b_n)_n$ positive Folgen reeller Zahlen. Zudem gilt gleichzeitig $a_n\in O(b_n)$ sowie $b_n\in O(a_n)$ für $n\to\infty$.
	\begin{parts}
		\part[2]
		\textbf{Behauptung:}
		\begin{center}
			Sowohl $(a_n)_n$ als auch $(b_n)_n$ sind beschränkte Folgen.
		\end{center}
		\begin{solution}
			falsch, $a_n = b_n = n$ sind unbeschränkte folgen und es gilt aber $n\in O(n)$
		\end{solution}
		\vspace{0.5cm}
		Beweise oder widerlege die Behauptung oder finde ein Gegenbeispiel. Begründe deine Antwort!
		\part[2]
		\textbf{Behauptung:}
		\begin{center}
			Es existiert eine Konstante $C>0$ mit $C^{-1} < a_n / b_n < C$ für alle $n\in\N$.
		\end{center}
		\vspace{0.5cm}
		Beweise oder widerlege die Behauptung oder finde ein Gegenbeispiel. Begründe deine Antwort!
		\begin{solution}
			Die Behauptung ist wahr. Da $a_n, b_n > 0$ für alle $n \in \mathbb{N}$, ist das Verhältnis $\frac{a_n}{b_n}$ stets wohldefiniert und positiv.

\begin{enumerate}
    \item \textbf{Anwendung der Definitionen:}
    \begin{itemize}
        \item Aus $a_n \in O(b_n)$ folgt: Es existieren $C_0 > 0$ und $N_0 \in \mathbb{N}$, sodass für alle $n \geq N_0$ gilt:
        \begin{equation} a_n \leq C_0 b_n \implies \frac{a_n}{b_n} \leq C_0 \end{equation}
        \item Aus $b_n \in O(a_n)$ folgt: Es existieren $C_1 > 0$ und $N_1 \in \mathbb{N}$, sodass für alle $n \geq N_1$ gilt:
        \begin{equation} b_n \leq C_1 a_n \implies \frac{1}{C_1} \leq \frac{a_n}{b_n} \end{equation}
    \end{itemize}

    \item \textbf{Zusammenführung für fast alle (alle ausser endlich viele) $n$:}\\
    Setze $N := \max\lrc{N_0, N_1}$. Für alle $n \geq N$ gilt:
    \begin{equation} \frac{1}{C_1} \leq \frac{a_n}{b_n} \leq C_0 \end{equation}

    \item \textbf{Erweiterung auf alle $n \in \mathbb{N}$:}\\
    Um die Schranken für die gesamte Folge zu finden, betrachten wir die endliche Menge der Quotienten vor dem Index $N$:
    \[ S = \left\{ \frac{a_1}{b_1}, \frac{a_2}{b_2}, \dots, \frac{a_{N-1}}{b_{N-1}} \right\} \]
    Wir definieren nun das globale Maximum $M$ und Minimum $m$:
    \begin{align*}
        M &:= \max\lr{S \cup \{C_0\}} \\
        m &:= \min\lr{S \cup \{1/C_1\}}
    \end{align*}
    Da alle beteiligten Werte positiv sind, gilt für alle $n \in \mathbb{N}$: $0 < m \leq \frac{a_n}{b_n} \leq M < \infty$.

    \item \textbf{Wahl der Konstante $C$:}\\
    Wähle $C := \max\lrc{ M + 1, \frac{1}{m} + 1}$. Da $M+1 > M$ und $\frac{1}{m}+1 > \frac{1}{m}$ (woraus $m > \frac{1}{M+1}$ folgt), ergibt sich für alle $n \in \mathbb{N}$:
    \[ C^{-1} < m \leq \frac{a_n}{b_n} \leq M < C \]
\end{enumerate}
Damit ist die Existenz einer solchen Konstante $C$ bewiesen.
		\end{solution}
		\droptotalpoints
	\end{parts}

	\question
	Begründe oder widerlege:
	\begin{parts}
		\part[1] $\lr{2n\sqrt{n}+10}\in O \lr{n^2}$
		\begin{solution}
			Wegen
			\begin{align*}
				 & \limsup_{n\to\infty}\frac{\abs{2n\sqrt{n}+10}}{\abs{n^2}} =                              \\
				 & =\limsup_{n\to\infty}\frac{\abs{n^2\lr{\frac{2}{\sqrt{n}}+\frac{10}{n^2}}}}{\abs{n^2}} = \\
				 & =\limsup_{n\to\infty}\lr{\frac{2}{\sqrt{n}}+\frac{10}{n^2}} =                            \\
				 & =\limsup_{n\to\infty}\lr{\frac{2}{\sqrt{n}}} + \limsup_{n\to\infty}\lr{\frac{10}{n^2}} = \\
				 & =0 < \infty,
			\end{align*}
			ist die Behauptung wahr.
		\end{solution}

		\part[1] $\lr{5\cos\lr{\pi m^2}m}\in O\lr{m}$
		\begin{solution}
			Wegen
			\begin{align*}
				 & \limsup_{m\to\infty}\frac{\abs{5\cos\lr{\pi m^2}m}}{\abs{m}} = \\
				 & =\limsup_{m\to\infty} 5\abs{\cos\lr{\pi m^2}} =                \\
				 & = 5 < \infty,
			\end{align*}
			ist die Behauptung wahr.
		\end{solution}

		\part[1] $2^{n+1}\in\ O\lr{2^n}$
		\begin{solution}
			Wegen
			\begin{align*}
				 & \limsup_{n\to\infty}\frac{\abs{2^{n+1}}}{\abs{2^n}} = \\
				 & =\limsup_{n\to\infty} 2 =                             \\
				 & =2 < \infty,
			\end{align*}
			ist die Behauptung wahr.
		\end{solution}

		\part[1] $2^{2n}\in O\lr{2^n}$
		\begin{solution}
			Es gilt
			\begin{align*}
				 & \limsup_{n\to\infty}\frac{\abs{2^{2n}}}{\abs{2^n}} =        \\
				 & =\limsup_{n\to\infty}\frac{\abs{2^n\cdot 2^n}}{\abs{2^n}} = \\
				 & =\limsup_{n\to\infty} = \infty.
			\end{align*}
			Deshalb ist die Behauptung falsch.
		\end{solution}
		\droptotalpoints
	\end{parts}


	\question[2] Erkläre genau, wie man fast jeden Algorithmus $A$  modifizieren kann, sodass er eine \enquote{gute} best-case Laufzeit hat.
	\begin{solution}
Aus didaktischen gründen beschreiben die Modifikation stets anhand von zwei verschiedenen Algorithmen: Dem $\Theta(n^2)$ \lstinline|bubble_sort| Algorithmus (siehe \cref{listing:bubble_pre}) zum Sortieren eines Arrays $L$ der Länge $n$ und einem $\Theta(n^3)$ Algorithmus \lstinline|det| (siehe \cref{listing:det_pre_check}) zur Berechnung des Determinanten einer gegeben $n\times n$-Matrix $M$. Die Modifikation für den allgmeinen Fall ist analog durchzuführen.

\begin{enumerate}
\item Prüfe, ob die erhaltene Eingabe eine besonders einfachen Instanzen ist:
\begin{itemize}
	\item Für \lstinline|bubble_sort|: Überprüfe, ob das Array $L$ bereits sortiert ist.\footnote{Eine solche Überprüfung (pre-Check) kann meist in kurzer Zeit durchgeführt werden (z.B. in $\Theta(n)$, um die Ordnung in einem Array der Länge $n$ zu überprüfen bzw. in $O(n^2)$, um die Struktur einer Matrix zu überprüfen).}
	\item Für \lstinline|det|: Überprüfe ob $M$ entweder eine Diagonalmatrix ist oder (mindestens) eine Nullzeile / Nullspalte besitzt.
\end{itemize}
\item Falls \enquote{ja}, berechne die Lösung für diese einfache Instanz (das Array $L$ bleibt unverändert, Determinate in diesen einfachen Fällen berechnen) und akzeptiere diese Lösung als Ausgabe.
\item Falls \enquote{nein}, führe den ursprünglichen Algorithmus $A$ aus und übernimm dessen Ausgabe.
\end{enumerate}

Natürlich wird durch diese Modifikation $A$ seine Laufzeit im worst-case verlangsamt (wegen der Kosten den pre-Check). Dennoch hat der modifizierte Algorithmus eine gute best-case Laufzeit, da er für einfache Instanzen sehr schnell die Lösung berechnen kann.

\lstinputlisting[basicstyle=\footnotesize\ttfamily,language=Python,caption=\texttt{bubble\_pre.py},label=listing:bubble_pre,numbers=left]{EF/Complexity/Code/bubble_pre.py}
\lstinputlisting[basicstyle=\footnotesize\ttfamily,language=Python,caption=\texttt{det\_pre\_check.py},label=listing:det_pre_check,numbers=left]{EF/Complexity/Code/det_pre_check.py}
\end{solution}


	\clearpage
	\question
	\begin{parts}
		\part[1] Jeder denkbare Algorithmus, der die Matrix-Matrix-Multiplikation $C = AB$ einer beliebigen $n\times n$ Matrix $A$ mit einer beliebigen $n\times n$ Matrix $B$ berechnet, hat eine Laufzeit von
		\begin{checkboxes}
			\choice $\Omega\lr{n^3}$
			\choice $\Theta\lr{n^3}$
			\choice $O\lr{n^3}$
			\correctchoice keine der obigen Antworten ist richtig
		\end{checkboxes}
		\begin{solution}
			Wir wissen, dass es Algorithmen gibt, welche eine gewünschte Matrix-Matrix-Multiplikation in $O\lr{n^{2.38}}$ implementieren. Damit ist die erste Aussage sicherlich falsch.
			Der Algorithmus
			\begin{lstlisting}[language=Python, numbers=none]
				# Multiply n x n matrices a and b
				def mmm(A, B, C, n):
					for i in range(n):
						for j in range(n):
							for k in range(n):
								C[i * n + j] += A[i * n + k] * B[k * n + j]
								for h in range(n):
									print(h)
			\end{lstlisting}
			implementiert eine gewünschte Matrix-Matrix-Multiplikation in Zeit $\Theta\lr{n^4}$. Somit sind auch die zweite und dritte Behauptung falsch.
		\end{solution}
		\part[1] Beschreibe ein algorithmisches Problem und einen Algorithmus für dieses Problem, welcher eine Laufzeit von $\Theta\lr{n}$ hat.
		\droptotalpoints
		\begin{solution}
			\begin{lstlisting}[language=Python, numbers=none]
				# finde maximales Element in Array A der Länge n
				def find_max(A):
					current_max = A[0]
					for element in A:
						if element > current_max:
							current_max = element
					return current_max
			\end{lstlisting}
		\end{solution}
	\end{parts}


	\question Gegeben sei der folgende Suchalgorithmus:
		\begin{lstlisting}[language=Python]
		def search(E, n, key):
			left = 0
			right = n - 1

			while left <= right:
				if E[left] == key or E[right] == key:
					return True

				left += 1
				right -= 1

			return False
		\end{lstlisting}
	Als Eingabe erhält er ein Array \verb|E| mit \verb|n| Einträgen sowie einen Schlüssel \verb|key| für den er überprüft, ob er in dem
	Eingabe-Array enthalten ist oder nicht.
	\begin{parts}
		\part[1] Bestimme die exakte Anzahl der Vergleiche im best-case.
		\begin{solution}
			Für $n = 0$ ist es eine Operation. Für $n\geq 1$ sind es jeweils zwei Operationen.
		\end{solution}
		\part[1] Die Laufzeit des Algorithmus im worst-case ist:
		\begin{checkboxes}
			\choice $\Theta\lr{n\log(n)}$
			\correctchoice $\Theta\lr{n}$
			\choice $\Theta\lr{n^2}$
			\choice $\Theta\lr{n^3}$
			\choice keine der obigen Antworten ist richtig
		\end{checkboxes}
	\end{parts} \droptotalpoints

	\clearpage
	\question Bestimme jeweils die Laufzeit der folgenden Algorithmen in Abhängigkeit von $n\in\N$.
	\begin{parts}
		\part[1] \leavevmode
		\begin{lstlisting}[language=Python]
			import random

			a = 0
			b = 2

			for i in range(n):
				for j in range(n):
					b += random.randint(0, RAND_MAX)
					a = a + b * random.randint(0, RAND_MAX)
			\end{lstlisting}
		\begin{checkboxes}
			\choice $\Theta\lr{n}$
			\choice $\Theta\lr{\sqrt{n}}$
			\choice $\Theta\lr{\log(n)}$
			\correctchoice $\Theta\lr{n^2}$
			\choice $\Theta\lr{n\log(n)}$
			\choice keine der obigen Antworten ist richtig
		\end{checkboxes}

		\part[1] \leavevmode
		\begin{lstlisting}[language=Python]
			import random

			maximum = 1000

			a = 0
			b = 0

			for i in range(n):
				a = a + random.randint(0, maximum)

			for j in range(n):
				b = b + random.randint(0, maximum)
		\end{lstlisting}
		\begin{checkboxes}
			\correctchoice $\Theta\lr{n}$
			\choice $\Theta\lr{\sqrt{n}}$
			\choice $\Theta\lr{\log(n)}$
			\choice $\Theta\lr{n^2}$
			\choice $\Theta\lr{n\log(n)}$
			\choice keine der obigen Antworten ist richtig
		\end{checkboxes}
		\part[1] \leavevmode
		\begin{lstlisting}[language=Python]
			a = 0
			i = n

			while i > 0:
				a += i
				i //= 2  # Integer division
		\end{lstlisting}
		\begin{checkboxes}
			\choice $\Theta\lr{n}$
			\choice $\Theta\lr{\sqrt{n}}$
			\correctchoice $\Theta\lr{\log(n)}$
			\choice $\Theta\lr{n^2}$
			\choice $\Theta\lr{n\log(n)}$
			\choice keine der obigen Antworten ist richtig
		\end{checkboxes}

		\part[1] \leavevmode
		\begin{lstlisting}[language=Python]
			i = n // 2
			k = 0

			while i <= n:
				j = 2
				while j <= n:
					k = k + n // 2
					j = j * 2
				i += 1
		\end{lstlisting}
		\begin{checkboxes}
			\choice $\Theta\lr{n}$
			\choice $\Theta\lr{\sqrt{n}}$
			\choice $\Theta\lr{\log(n)}$
			\choice $\Theta\lr{n^2}$
			\correctchoice $\Theta\lr{n\log(n)}$
			\choice keine der obigen Antworten ist richtig
		\end{checkboxes}
	\end{parts} \droptotalpoints

	\question Der Algorithmus \lstinline|bubble_sort| ist gegeben durch:
	\begin{lstlisting}[language=Python]
		# A function to implement bubble sort
		def bubble_sort(arr):
			n = len(arr)
			for i in range(n - 1):
				# Last i elements are already in place
				for j in range(n - i - 1):
					if arr[j] > arr[j + 1]:
						# Swap elements
						arr[j], arr[j + 1] = arr[j + 1], arr[j]
	\end{lstlisting}
	\begin{parts}
		\part[2] Bestimme die worst-case Laufzeit von \lstinline|bubble_sort|. Zeige deine Schritte ausführlich.
		\begin{solution}
			Die Anzahl der durchgeführten Vergleiche ist
			\begin{align*}
				(n-1) + (n-2) + (n-3) +\ldots + 2 + 1 = n(n-1)/2.
			\end{align*}
		\end{solution}
		\part[2] Bestimme ebenso die best-case Laufzeit von \lstinline|bubble_sort|. Erkläre deine Überlegungen. \droptotalpoints
		\begin{solution}
			Im best-case ist das Array schon sortiert und die Vergleiche stets falsch. Dennoch ist die Laufzeit $\Theta\lr{n^2}$, da die Schleifen ohnehin ausgeführt werden.
		\end{solution}
	\end{parts}
	\clearpage
	\noindent
	Es seien $f$ und $g$ Folgen reeller Zahlen.\\
	\vspace{0.5cm}
	\centering
	\small
	\begin{tabular}{|c|c|c|}
		\hline
		\textbf{Notation} & \textbf{Definition}                                              & \textbf{Bedeutung}                                  \\
		\hline
		$f \in o(g)$      & $\lim\limits_{n \to \infty} \abs{\frac{f(n)}{g(n)}} = 0$         & $f$ wächst langsamer als $g$                        \\
		\hline
		$f \in O(g)$      & $\limsup\limits_{n \to \infty} \abs{\frac{f(n)}{g(n)}} < \infty$ & $f$ wächst höchstens so schnell wie $g$             \\
		\hline
		$f \in \Omega(g)$ & $\liminf\limits_{n \to \infty} \abs{\frac{f(n)}{g(n)}} > 0$      & \makecell{$f$ wächst mindestens so schnell wie $g$, \\$g\in O(f)$}                          \\
		\hline
		$f \in \Theta(g)$ & \makecell{$f \in O(g)\cap \Omega(g)$,                                                                                  \\ $0 < \liminf\limits_{n \to \infty} \abs{\frac{f(n)}{g(n)}} \leq$ \\ $\leq\limsup\limits_{n \to \infty} \abs{\frac{f(n)}{g(n)}} < \infty$} & \makecell{$f$ wächst gleich schnell wie $g$,        \\sowohl $f\in O(g)$ als auch $f\in\Omega(g)$} \\
		\hline
		$f \in \omega(g)$ & $\lim\limits_{n \to \infty} \abs{\frac{f(n)}{g(n)}} = \infty$    & \makecell{$f$ wächst schneller als $g$,             \\$g\in o(f)$}                                      \\
		\hline
	\end{tabular}

	(Wir werden meist nur $\lim$ anstelle von $\limsup$ oder $\liminf$ betrachten müssen.)
\end{questions}
